\documentclass[../../main.tex]{subfiles}

\begin{document}

\section{Modules and Libraries}

First, what are modules and libraries?

\begin{definition}{}{}
\textbf{Modules} in Python are files that consists of definitions
and statements for organization and reusability. \textbf{Libraries}
are collection of modules and pre-written packages that can later
be imported and distributed.
\end{definition}

One of the reason why Python is so popular in ML and data science
is because of the extensive libraries. Having the pre-written
libraries specifically for ML and tensors like PyTorch simply
makes our lives easier.

\subsection{NumPy}

NumPy is a Python package that helps handle multidimensional arrays.
Let's get started by importing NumPy to our code. We could import
it as whatever alias we want, but \verb|np| is the standard for ML,
so let's stick with that.

\begin{jupyter}
\begin{lstlisting}[style=jin]
import numpy as np
\end{lstlisting}
\begin{lstlisting}[style=jout]
\end{lstlisting}
\end{jupyter}

Assuming that you have read the previous note or have strong
foundations in linear algebra, we can implement the mathematical
properties in Python with NumPy. Below is an example of how we can
define vectors in different dimensions. Note that \verb|.ndim| can
be used to find the dimension of a vector.

\begin{jupyter}
\begin{lstlisting}[style=jin]
arr_1d = np.array([1, 2, 3])
arr_2d = np.array([[1, 2, 3],
                   [4, 5, 6]])
arr_3d = np.array([[[1, 2], [3, 4]],
                   [[5, 6], [7, 8]],
                   [[9, 10], [11, 12]]])

print(f"Dimension: {arr_1d.ndim}, Shape: {arr_1d.shape}")
print(f"Dimension: {arr_2d.ndim}, Shape: {arr_2d.shape}")
print(f"Dimension: {arr_3d.ndim}, Shape: {arr_3d.shape}")
\end{lstlisting}

\begin{lstlisting}[style=jout]
Dimension: 1, Shape: (3,)
Dimension: 2, Shape: (2, 3)
Dimension: 3, Shape: (3, 2, 2)
\end{lstlisting}
\end{jupyter}

A simple way to understand multidimensional array in NumPy is
stacking previous layers. For instance, \verb|arr_2d| has the
shape $(2, 3)$. We can think of it as stacking two 1D array
with three elements. Similarly, \verb|arr_3d| has the shape
$(3, 2, 2)$. We could think of it as stacking three layers of
2D arrays that consists of two 1D arrays. It is quite hard
to visualize higher dimensional arrays, so let's take a look
at different things that NumPy can do.

Below is an example where we utilize arange, reshape, dtype,
and slicing techniques in NumPy. Let's create an array first.

\begin{jupyter}
\begin{lstlisting}[style=jin]
# Making an array with numbers 0-17 and
# rerranging it to shape (3, 3, 2)
a = np.arange(18).reshape(3, 3, 2)
a
\end{lstlisting}

\begin{lstlisting}[style=jout]
array([[[ 0,  1],
        [ 2,  3],
        [ 4,  5]],

       [[ 6,  7],
        [ 8,  9],
        [10, 11]],

       [[12, 13],
        [14, 15],
        [16, 17]]])
\end{lstlisting}
\end{jupyter}

We can also find the data type with the attribute \verb|.dtype|.

\begin{jupyter}
\begin{lstlisting}[style=jin]
a.dtype.name
\end{lstlisting}

\begin{lstlisting}[style=jout]
'int64'
\end{lstlisting}
\end{jupyter}

This means that the individual elements in array \verb|a| are 8
bytes integers. Just as we index and slice 1D array, we could
do the same with NumPy. Consider the following example for a
matrix.

\begin{jupyter}
\begin{lstlisting}[style=jin]
a = np.arange(9).reshape(3, 3)
print(f"{a}\n")

# First row, second column element
print(a[0, 1])
# Entire second row
print(a[1, :])
# Entire third column
print(a[:, 2])
\end{lstlisting}

\begin{lstlisting}[style=jout]
[[0 1 2]
 [3 4 5]
 [6 7 8]]

1
[3 4 5]
[2 5 8]
\end{lstlisting}
\end{jupyter}

This was the basic introduction to NumPy. Please check out the
\href{https://numpy.org/doc/stable/}{documentation} for further
information.

\subsection{Matplotlib}

Do you want to visualize something with Python? Matplotlib is here
for your need! Matplotlib is a Python library for visualization.
Let's first import the library to take a look at an example.

\begin{jupyter}
\begin{lstlisting}[style=jin]
import matplotlib.pyplot as plt
%matplotlib inline
\end{lstlisting}
\begin{lstlisting}[style=jout]
\end{lstlisting}
\end{jupyter}

If you are using Jupyter Notebook like the examples in these
notes, it is highly recommended to keep \verb|%matplotlib inline|
to make our lives easier by plotting directly in the notebook
without an external window. Although Matplotlib is an essential
library for visualizing model performance and analyzing, let's
stick with basic plotting for now as I believe more information
would be an overkill for this note.

For this brief introduction, let's plot Sigmoid function, which
we will discuss in later notes.

\begin{jupyter}
\begin{lstlisting}[style=jin]
# Define domain and the function
x = np.linspace(-10, 10, 100)
y = 1 / (1 + np.exp(-x))

plt.figure()
plt.plot(x, y, label='Sigmoid')
plt.grid(True)
plt.xlabel('x')
plt.ylabel('sigmoid(x)')
plt.title('Sigmoid Function')
plt.legend()
plt.show()
\end{lstlisting}
\vspace{-0.6\baselineskip}
\includegraphics[scale=0.43]{CS0102.26011-1}
\end{jupyter}

I know this intro was quite brief, but we will continue using
amazing features in Matplotlib! For more information, please
visit the official
\href{https://matplotlib.org/stable/tutorials/index}{tutorial page}.

In the next note, we will discuss foundational terms and concepts
in ML. It will also include a high level introduction to large
language model (LLM). Understanding the high level overview of how 
LLMs operate under the hood really helped me grasp the intuition.
Without further ado, let's get started!
\end{document}

